\section{Polytopes}
A polytope is the convex hull of finitely many points. Its minimal
representation consists exactly of its extreme points, or vertices.
For a linear functional, the value at a convex combination is the same
combination of the values at its generators. Hence its maximum is the
largest generator value; the full maximizing set is the convex hull of
\emph{all} generators tying for that largest value. The analogous statement
holds for minima. A $k$-face has dimension $k$; vertices have dimension zero,
edges dimension one, and facets of an $n$-polytope dimension $n-1$.

\practice{8.5.P1}{Turn five inequalities into a vertex list}
In the first quadrant the constraints are
$x+3y\le12$, $x+2y\le9$, $2x+y\le12$.
The axis vertices are $(0,0),(6,0),(0,4)$, obtained by imposing the most
restrictive intercept on each axis. Pairwise boundary intersections are
$(3,3),(5,2),(24/5,12/5)$. The last violates $x+2y\le9$, because
$24/5+24/5=48/5>9$. The first two pass the remaining inequalities.
Thus the minimal representation is
\[
 P=\operatorname{conv}\{(0,0),(6,0),(5,2),(3,3),(0,4)\}.
\]
Each retained point is a genuine corner of the feasible polygon; no retained
point is a convex combination of the others. The printed final sentence
calls these ``three'' vertices while listing five; there are five.

\exercise{8.5.1}{All maximizers of three functionals on a triangle}
At $p_1=(1,0),p_2=(2,3),p_3=(-1,2)$, the functional values are
\[
 \begin{array}{c|rrr|c|l}
 &p_1&p_2&p_3&\max&\text{all maximizing points}\\\hline
 x-y&1&-1&-3&1&p_1\\
 x+y&1&5&1&5&p_2\\
 -3x+y&-3&-3&5&5&p_3
 \end{array}
\]
These respectively answer (a), (b), (c). To justify completeness, write
$x=\sum_i\lambda_ip_i$ with nonnegative coefficients of sum one.
If a positive coefficient were attached to a vertex with value below the
maximum, their weighted average would be strictly below that maximum.
Here each maximum occurs at just one vertex, so no other point can attain it.

\exercise{8.5.2}{Two maximizing edges}
For $p_1=(0,-1),p_2=(2,1),p_3=(1,2)$, evaluation gives
\[
 \begin{array}{c|rrr|c|l}
 &p_1&p_2&p_3&\max&\text{all maximizing points}\\\hline
 x+y&-1&3&3&3&\operatorname{conv}\{p_2,p_3\}\\
 x-y&1&1&-1&1&\operatorname{conv}\{p_1,p_2\}\\
 -2x+y&-1&-3&0&0&p_3
 \end{array}
\]
In each tied case every point on the listed edge has the same value by
linearity. A convex combination involving the lower-value vertex with
positive weight has smaller value, so the whole maximizing set is exactly that edge.

\exercise{8.5.3}{Minima for the first triangle}
Use the values from Exercise 8.5.1 and select the smallest in each row.
For (a), the minimum of $x-y$ is $-3$, attained only at $p_3$.
For (b), the minimum of $x+y$ is $1$, attained at every point of the edge
$\operatorname{conv}\{p_1,p_3\}$ and nowhere else.
For (c), the minimum of $-3x+y$ is $-3$, attained exactly on
$\operatorname{conv}\{p_1,p_2\}$.
The weighted-average argument used for maxima also proves these sets contain all minimizers.

\exercise{8.5.4}{Minima for the second triangle}
From the values in Exercise 8.5.2:
(a) $x+y$ has minimum $-1$ only at $p_1$;
(b) $x-y$ has minimum $-1$ only at $p_3$;
(c) $-2x+y$ has minimum $-3$ only at $p_2$.
No other point can attain one of these strict vertex minima, because a
combination using another vertex would increase the weighted average.

\exercise{8.5.5}{A polygon bounded by two sloping lines}
The inequalities are $x+2y\le10$, $3x+y\le15$, $x,y\ge0$.
The axis endpoints are $(5,0)$ and $(0,5)$. Solve both boundary equations:
substitute $y=15-3x$ into the first, obtaining $x+30-6x=10$,
so $x=4,y=3$. This intersection is in the first quadrant.
Together with the origin these are all corners:
\[
 P=\operatorname{conv}\{(0,0),(5,0),(4,3),(0,5)\}.
\]
Each corner is the meeting point of two boundary lines of the feasible
polygon, and each lies on its boundary with a change in direction, so none
is redundant in the minimal representation.

\exercise{8.5.6}{Another two-constraint polygon}
Here $2x+3y\le18$, $4x+y\le16$, $x,y\ge0$.
The axis endpoints are $(4,0),(0,6)$.
The two sloping boundaries intersect where $y=16-4x$ and
$2x+3(16-4x)=18$, giving $x=3,y=4$.
Thus
\[
 P=\operatorname{conv}\{(0,0),(4,0),(3,4),(0,6)\}.
\]
All listed points satisfy all constraints and are the successive corners
of the feasible polygon, hence give its minimal representation.

\exercise{8.5.7}{Three boundaries meet at one corner}
The constraints are $x+3y\le18$, $x+y\le10$, $4x+y\le28$,
with nonnegative coordinates. The restrictive axis intercepts are $(7,0)$
and $(0,6)$. Solving any pair of the sloping equalities gives the same
point $(6,4)$, which satisfies all three with equality.
Hence there is only one non-axis corner, not three:
\[
 P=\operatorname{conv}\{(0,0),(7,0),(6,4),(0,6)\}.
\]
The inequality $x+y\le10$ touches this corner but introduces no new vertex.

\exercise{8.5.8}{A redundant inequality}
The constraints are $2x+y\le8$, $x+y\le6$, $x+2y\le7$,
with $x,y\ge0$. Their restrictive axis endpoints are $(4,0),(0,7/2)$.
The first and third boundaries meet at $(3,2)$, which satisfies $x+y\le6$.
The intersections with the middle boundary are $(2,4)$ and $(5,1)$;
they violate respectively $x+2y\le7$ and $2x+y\le8$ and must be discarded.
Therefore
\[
 P=\operatorname{conv}\{(0,0),(4,0),(3,2),(0,7/2)\}.
\]
The discarded points are intersections of boundary lines, but they are
not feasible and therefore cannot be vertices of this polytope.

\exercise{8.5.9}{An extreme point that is not an exposed vertex}
Let $D=\{(x,y):x^2+(y-1)^2\le1\}$ and
$K=\operatorname{conv}(D\cup\{(3,0)\})$.
Every point of $K$ has $y\ge0$, and its points with $y=0$ form precisely
the segment $[(0,0),(3,0)]$. Indeed the disk itself touches $y=0$ only
at the origin.
If the origin were a strict convex combination of two distinct points in
$K$, their nonnegative heights would both have to be zero. Their
nonnegative horizontal coordinates on that segment could then average to
zero only if both points were the origin. Thus the origin is extreme.
However, the only supporting line to the disk at the origin is its tangent
$y=0$. Any supporting line to $K$ there must also support the contained
disk, so it must be that same line. It intersects $K$ in the whole segment
to $(3,0)$, not in a singleton.
\answer{The origin is extreme but is not a vertex in the section's
supporting-hyperplane (zero-dimensional exposed-face) sense.}

\exercise{8.5.10}{A closed unbounded counterexample to profile generation}
Take the ray $S=\{(x,0):x\ge0\}$. It is closed and convex.
Its origin is extreme, while every point with $x>0$ is the midpoint of
two other ray points. Thus its profile is $\{(0,0)\}$.
The convex hull of that profile is just the singleton, not the whole ray.
This shows why closedness without boundedness is insufficient.

\exercise{8.5.11}{A bounded but nonclosed counterexample}
Take the half-open segment $S=\{(x,0):0\le x<1\}$.
It is bounded and convex. Its profile is again just the origin, since
any other point has two distinct nearby points on either side inside the segment.
The hull of the nonempty profile is a singleton, not $S$.
This time the missing hypothesis is closedness.

\exercise{8.5.12}{Face counts for simplices}
\step{(a) Count the faces of the five-simplex.}
A five-simplex has six vertices; each $k$-face is determined by choosing
$k+1$ of them. Thus its counts for $k=0,1,2,3,4$ are
$6,15,20,15,6$. Their alternating sum is $6-15+20-15+6=2$,
equal to $1+(-1)^4$, as Euler's formula requires.
\step{(b) Display the pattern and explain it.}
\[
 \begin{array}{c|rrrrr}
 n\backslash k&0&1&2&3&4\\\hline
 1&2&-&-&-&-\\2&3&3&-&-&-\\3&4&6&4&-&-\\
 4&5&10&10&5&-\\5&6&15&20&15&6
 \end{array}
\]
The general formula is $f_k(S^n)=\binom{n+1}{k+1}$ for $0\le k\le n-1$.
Affine independence ensures that each chosen $k+1$ vertices defines a
$k$-simplex face and that different subsets give different faces.
Dashes denote dimensions not belonging to the list of proper faces.

\exercise{8.5.13}{Face counts for hypercubes}
\step{(a) Count a five-cube and check Euler.}
Choose $k$ coordinates allowed to vary and fix each remaining coordinate
at either endpoint. Thus $f_k(C^n)=2^{n-k}\binom nk$.
For $n=5$ the counts are $32,80,80,40,10$;
$32-80+80-40+10=2=1+(-1)^4$.
\step{(b) The requested chart.}
\[
 \begin{array}{c|rrrrr}
 n\backslash k&0&1&2&3&4\\\hline
 1&2&-&-&-&-\\2&4&4&-&-&-\\3&8&12&6&-&-\\
 4&16&32&24&8&-\\5&32&80&80&40&10
 \end{array}
\]
The coordinate-choice argument proves the formula, rather than only
predicting it from these few values.

\exercise{8.5.14}{Crosspolytopes in increasing dimensions}
\step{(a) Draw the first two cases.}
$X^1$ is the segment joining $-v_1$ to $v_1$.
$X^2$ is the parallelogram with successive vertices
$v_1,v_2,-v_1,-v_2$. Below are representatives using standard-coordinate
generators; an invertible linear change gives the analogous pictures for
any independent generators.
\begin{center}\setlength{\unitlength}{1mm}
\begin{picture}(125,40)
\put(5,20){\line(1,0){40}}\put(5,20){\circle*{1.3}}\put(45,20){\circle*{1.3}}
\put(1,14){$-v_1$}\put(42,14){$v_1$}\put(22,25){$X^1$}
\put(85,20){\polyline(18,0)(0,16)(-18,0)(0,-16)(18,0)}
\put(106,19){$v_1$}\put(83,37){$v_2$}\put(56,19){$-v_1$}\put(82,0){$-v_2$}
\end{picture}\end{center}
\step{(b) The three-dimensional case.}
There are $6$ vertices, $12$ edges, and $8$ triangular facets.
The standard-coordinate version is a regular octahedron; for arbitrary
independent generators it is a linearly distorted octahedron with the same face counts.
\step{(c) The four-dimensional case.}
Its counts are $8,24,32,16$. The alternating sum is
$8-24+32-16=0=1+(-1)^3$, as required.
\step{(d) Explain the general formula.}
A proper face selects $k+1$ of the $n$ opposite vertex pairs, and selects
one sign from each selected pair; both opposite vertices cannot belong to
one proper face. These choices give
\[
 f_k(X^n)=2^{k+1}\binom n{k+1},\qquad0\le k\le n-1.
\]
Every such selection is a simplex face: for the standard crosspolytope a
supporting functional assigning the selected signed coordinates value one
and all other coordinate coefficients zero exposes exactly that face.
Linear equivalence preserves the face structure for general independent generators.

\exercise{8.5.15}{Face counts for a pyramid}
Let the $(n-1)$-dimensional base be $Q$ and the apex lie outside its affine hull.
\step{(a) Vertices.}
All base vertices remain vertices and the apex adds one, so
$f_0(P^n)=f_0(Q)+1$.
\step{(b) Intermediate dimensions.}
A $k$-face not containing the apex is a $k$-face of $Q$.
A $k$-face containing the apex is the pyramid over a $(k-1)$-face of $Q$.
Thus $f_k(P^n)=f_k(Q)+f_{k-1}(Q)$ for $1\le k\le n-2$.
\step{(c) Facets.}
There is one base facet $Q$; each of its $(n-2)$-faces joined to the apex
gives a side facet. Therefore $f_{n-1}(P^n)=1+f_{n-2}(Q)$, for $n\ge2$.
These types are disjoint and cover every face of a pyramid.

\exercise{8.5.16}{Four statements about profiles and simplices}
\begin{enumerate}[label=(\alph*)]
\item \textbf{True.} This is the definition of a polytope.
\item \textbf{True.} An extreme point cannot lie strictly between two
other points of the convex set. If $p\in uv$ and $p\ne u$, it must equal $v$.
\item \textbf{False.} Compactness is necessary in the profile theorem;
the ray in Exercise 8.5.10 is a counterexample.
\item \textbf{True.} A four-simplex has five vertices; omit one at a
time to obtain its five facets, each a three-simplex, or tetrahedron.
\end{enumerate}

\exercise{8.5.17}{Four more statements about faces}
\begin{enumerate}[label=(\alph*)]
\item \textbf{False.} A three-dimensional cube has six square facets.
\item \textbf{True.} For polytopes the minimal-representation theorem
identifies extreme points with vertices. The distinction seen for a nonpolytope
in Exercise 8.5.9 does not occur for polytopes.
\item \textbf{False.} A maximizing point need not itself be extreme.
On the square $[-1,1]^2$, $f(x,y)=x$ is maximal at $(1,0)$, the midpoint
of the edge joining $(1,-1)$ and $(1,1)$.
\item \textbf{True.} A two-dimensional polytope is a polygon; each edge
joins successive vertices around its boundary, giving equally many vertices and edges.
\end{enumerate}

\exercise{8.5.18}{Removing an extreme point preserves convexity}
Suppose $v$ is extreme and $x,y\in S\setminus\{v\}$.
Their segment lies in $S$ by convexity and cannot pass through $v$ in its
interior, by extremality; neither endpoint is $v$. Thus the whole segment
lies in $S\setminus\{v\}$, proving that set convex.
Conversely, if $S\setminus\{v\}$ is convex and $v$ were strictly between
some distinct $x,y\in S$, neither could equal $v$.
Their segment would lie in the set with $v$ removed, contradicting its
containing $v$. Hence $v$ is extreme. Singleton cases work because the empty set is convex.

\exercise{8.5.19}{Add two positive dilations of one convex set}
An arbitrary element of $cS+dS$ is $cx+dy$ with $x,y\in S$.
Since $c,d>0$,
\[
 cx+dy=(c+d)\left(\frac c{c+d}x+\frac d{c+d}y\right)\in(c+d)S
\]
by convexity. This proves one inclusion. Conversely, for each $s\in S$,
$(c+d)s=cs+ds\in cS+dS$, proving the reverse inclusion.
Thus $cS+dS=(c+d)S$.

\exercise{8.5.20}{Convexity cannot be omitted}
Take $S=\{(0,0),(1,0)\}$ and $c=d=1$.
Then $S+S=\{(0,0),(1,0),(2,0)\}$, whereas
$2S=\{(0,0),(2,0)\}$. The point $(1,0)$ occurs only in the sum.
Thus the equality from Exercise 8.5.19 can fail for a nonconvex set.

\exercise{8.5.21}{A Minkowski sum of convex sets is convex}
Take two points $x=a_1+b_1$ and $y=a_2+b_2$ in $A+B$.
For $0\le t\le1$,
\[
 (1-t)x+ty=((1-t)a_1+ta_2)+((1-t)b_1+tb_2).
\]
The first parenthesis belongs to $A$ and the second to $B$ by their
convexity. Thus their sum is in $A+B$. This proves the entire segment is
included, as required.

\exercise{8.5.22}{Why the list of regular polyhedra has five entries}
\step{(a) Count incidences in two ways.}
Every one of the $r$ facets has $k$ edges, giving $kr$ facet-edge
incidences. Each edge belongs to two facets, giving $2e$; hence $kr=2e$.
Likewise every one of the $v$ vertices meets $s$ edges, while every edge
has two endpoints, so $sv=2e$.
\step{(b) Substitute those counts into Euler's formula.}
From $v-e+r=2$, substitute $v=2e/s$ and $r=2e/k$, then divide by $2e$:
\[
 \frac1s+\frac1k=\frac12+\frac1e.
\]
\step{(c) Exhaust the possible integer pairs.}
A face has at least three sides, so $k\ge3$; at least three edges meet at
any vertex, so $s\ge3$. The equation makes $1/s+1/k>1/2$.
If $k\ge6$, its left side is at most $1/3+1/6=1/2$, impossible.
For $k=3$, the strict inequality gives $s<6$, so $s=3,4,5$.
For $k=4$, it gives $s<4$, so $s=3$.
For $k=5$, it gives $s<10/3$, again $s=3$.
The counts follow from $e=(1/s+1/k-1/2)^{-1}$ and the incidence formulas:
\[
 \begin{array}{cc|rrr|l}
 k&s&v&e&r&\text{polyhedron}\\\hline
 3&3&4&6&4&\text{tetrahedron}\\
 4&3&8&12&6&\text{cube}\\
 3&4&6&12&8&\text{octahedron}\\
 5&3&20&30&12&\text{dodecahedron}\\
 3&5&12&30&20&\text{icosahedron}
 \end{array}
\]
These are the five realized regular convex polyhedra named in the exercise,
and the integer argument rules out every other possibility.
